% source: An algebraic approach to the Riemann-Roch theorem and the arithmetic theory of function fields (Athanasios Papaioannou), http://www.fen.bilkent.edu.tr/~franz/mat/Riemann-Roch.pdf
% pdf_page: 0010
% retrieved: 2026-07-26
% transcribed_by: page-transcriber (AI vision, no OCR)

\documentclass[11pt]{article}
\usepackage{amsmath,amssymb,amsthm}

% Small-caps theorem style matching the source.
\newtheoremstyle{scplain}{\topsep}{\topsep}{\itshape}{0pt}{\scshape}{.}{ }{\thmname{#1} \thmnote{#3}}
\theoremstyle{scplain}
\newtheorem*{definition}{Definition}
\newtheorem*{theorem}{Theorem}
\newtheorem*{lemma}{Lemma}

% Small-caps "Proof." to match the source.
\makeatletter
\renewenvironment{proof}[1][\proofname]{%
  \par\pushQED{\qed}\normalfont\topsep6\p@\@plus6\p@\relax
  \trivlist\item[\hskip\labelsep\scshape #1\@addpunct{.}]\ignorespaces}%
  {\popQED\endtrivlist\@endpefalse}
\makeatother
\renewcommand{\qedsymbol}{$\square$}

\DeclareMathOperator{\Div}{Div}

\begin{document}

\noindent
extend naturally to $A_{\mathbb{F}}$ by setting
$v_P(\alpha) = v_P(\alpha_P)$, where $\alpha_P$ is the $P$'th component of
$\alpha$. By definition, $v_P(\alpha_P) \geq 0$ for almost all
$P \in \mathbf{P}_{\mathbb{F}}$.

Adding to the sea of definitions there are already, we will define the following
$\mathbb{K}$-subspaces of $A_{\mathbb{F}}$:

\begin{definition}[1.12]
For $D \in \Div(\mathbb{F})$, we set
\[
  A_{\mathbb{F}}(D) =
    \{\alpha \in A_{\mathbb{F}} \mid v_P(\alpha) \geq -v_P(A)
      \text{ for all } P \in \mathbf{P}_{\mathbb{F}}\}.
\]
\end{definition}

Our first observation about adeles, and in fact a preliminary version of the
Riemann-Roch theorem, is the following:

\begin{theorem}[1.18]
For any divisor $D \in \Div(\mathbb{F})$,
\[
  \dim D = \deg D + 1 - g
    + \dim(A_{\mathbb{F}}/A_{\mathbb{F}}(D) + \mathbb{F}).
\]
\end{theorem}

\medskip\noindent{\scshape Proof.} Note that despite the infinite dimension of
$A_{\mathbb{F}}$, $A_{\mathbb{F}}(D)$ and $\mathbb{F}$, the above theorem
asserts that the quotient space
$A_{\mathbb{F}}/(A_{\mathbb{F}}(D)+\mathbb{F})$ is finite-dimensional over
$\mathbb{K}$.

We will complete the proof after three lemmas. We begin with

\begin{lemma}[1.19]
Let $D_1,D_2 \in \Div(\mathbb{F})$ be such that $D_1 \leq D_2$. Then
$A_{\mathbb{F}}(D_1) \subseteq A_{\mathbb{F}}(D_2)$ and
\[
  \dim(A_{\mathbb{F}}(D_2)/A_{\mathbb{F}}(D_1)) = \deg D_2 - \deg D_1.
\]
\end{lemma}

\begin{proof}
The first assertion of the lemma is obvious; it follows from the definition of
the spaces involved and the condition on $D_1,D_2$. For the second claim, note
that it is sufficient to prove it for in the case $D_2 = D_1 + P$, with
$P \in \mathbf{P}_{\mathbb{F}}$, since the general case then follows by an
inductive argument. Pick $t \in \mathbb{F}$ with
$v_P(t) = v_P(D_1) + 1$ and consider the $\mathbb{K}$-linear map
$\phi : A_{\mathbb{F}}(D_2) \longrightarrow \mathbb{F}_P$ defined by
$\alpha \mapsto (t\alpha)_P(P)$. We can easily see that $\phi$ is in fact
surjective, with kernel $A_{\mathbb{F}}(D_1)$. Therefore,
$\deg D_2 - \deg D_1 = \deg P = [\mathbb{F} : \mathbb{K}(x)]
= \dim(A_{\mathbb{F}}(D_2)/A_{\mathbb{F}}(D_1))$, and this implies the desired
result.
\end{proof}

We now go on to show the following

\begin{lemma}[1.20]
As before, let $D_1,D_2 \in \Div(\mathbb{F})$ be such that $D_1 \leq D_2$.
Then
\[
  \dim((A_{\mathbb{F}}(D_2)+\mathbb{F})/(A_{\mathbb{F}}(D_1)+\mathbb{F}))
    = (\deg D_2-\dim D_2) - (\deg D_1-\dim D_1).
\]
\end{lemma}

\begin{proof}
We have the following exact sequence of linear transformations
{\small
\[
  0 \longrightarrow L(D_2)L(D_1)
    \xrightarrow{\ \sigma_1\ } A_{\mathbb{F}}(D_2)/A_{\mathbb{F}}(D_1)
    \xrightarrow{\ \sigma_2\ }
    (A_{\mathbb{F}}(D_2)+\mathbb{F})/((A_{\mathbb{F}}(D_1)+\mathbb{F})
    \longrightarrow 0,
\]
}
where $\sigma_1$ and $\sigma_2$ are defined in the obvious manner. We further
claim that $\ker\sigma_2 \subseteq \operatorname{Im}\sigma_1$. Indeed, let
$\alpha \in A_{\mathbb{F}}(D_2)$ with
$\sigma_2(\alpha+A_{\mathbb{F}}(D_1))=0$. Then,
$\alpha \in A_{\mathbb{F}}(D_1)+\mathbb{F}$, so there is some
$x \in \mathbb{F}$ such that $\alpha-x \in A_{\mathbb{F}}(D_1)$. As
$A_{\mathbb{F}}(D_1) \subseteq A_{\mathbb{F}}(D_2)$ we conclude that
$x \in A_{\mathbb{F}}(D_2) \cap \mathbb{F} = L(D_2)$. Thus,
$\alpha+A_{\mathbb{F}}(D_1) = x+A_{\mathbb{F}}(D_1)
= \sigma_1(x+L(D_1))$ lies in the image of $\sigma_1$.

From the exactness of the previous sequence, we obtain that
{\small
\[
\begin{split}
  \dim((A_{\mathbb{F}}(D_2)+\mathbb{F})/(A_{\mathbb{F}}(D_1)+\mathbb{F}))
  &= \dim(A_{\mathbb{F}}(D_2)/A_{\mathbb{F}}(D_1))
     - \dim(L(D_2)/L(D_1)) \\
  &= (\deg D_2-\dim D_2) - (\deg D_1-\dim D_1),
\end{split}
\]
}
as desired.
\end{proof}

The third and last lemma can be stated as

\begin{lemma}[1.21]
If $B$ is a divisor such that $\dim B = \deg B + 1 - g$, then
$A_{\mathbb{F}} = A_{\mathbb{F}}(B) + \mathbb{F}$.
\end{lemma}

\begin{proof}
First, observe that for $B_1 \geq B$, we have
\[
  \dim B_1 \leq \deg B_1 + \dim B - \deg B = \deg B_1 + 1 - g,
\]
by lemma 1.13. On the other hand, $\dim B_1 \geq \deg B + 1 - g$ by Riemann's
inequality; therefore, $\dim B_1 = \deg B_1 + 1 - g$. Now, let
$\alpha \in A_{\mathbb{F}}$; by definition, we can find a divisor $B_1 > B$
such that $\alpha \in A_{\mathbb{F}}(B_1)$. By the proof of our previous lemma,
{\small
\[
\begin{aligned}
  &\dim((A_{\mathbb{F}}(B_1)+\mathbb{F})/(A_{\mathbb{F}}(B)+\mathbb{F})) \\
  &\qquad= (\deg B_1-\dim B_1) - (\deg B-\dim B) \\
  &\qquad= (g-1)-(g-1)=0.
\end{aligned}
\]
}
This yields the desired result, since it implies
$A_{\mathbb{F}}(B)+\mathbb{F} = A_{\mathbb{F}}(B_1)+\mathbb{F}$ and since
$\alpha \in A_{\mathbb{F}}(B_1)$, it follows that
$\alpha \in A_{\mathbb{F}}(B)+\mathbb{F}$, as desired.
\end{proof}

\begin{center}10\end{center}

\end{document}
